% p304_1; source PDF page 304; manually verified editable table.
\begin{tikzpicture}[x=1cm,y=-1cm,inner sep=0pt,outer sep=0pt]
\path[use as bounding box] (0,0) rectangle (16.000,14.445);
\draw[line width=.4pt] (0.000,0.000) rectangle (2.900,0.815);
\node[align=center,text width=2.640cm,font={\linespread{1}\fontsize{10}{13.50}\selectfont\bfseries}] at (1.450,0.407) {误解的类型};
\draw[line width=.4pt] (2.900,0.000) rectangle (7.300,0.815);
\node[align=center,text width=4.140cm,font={\linespread{1}\fontsize{10}{13.50}\selectfont\bfseries}] at (5.100,0.407) {具体情形};
\draw[line width=.4pt] (7.300,0.000) rectangle (16.000,0.815);
\node[align=center,text width=8.440cm,font={\linespread{1}\fontsize{10}{13.50}\selectfont\bfseries}] at (11.650,0.407) {范例};
\draw[line width=.4pt] (0.000,0.815) rectangle (2.900,14.445);
\node[align=left,text width=2.640cm,font={\linespread{1}\fontsize{10}{13.50}\selectfont}] at (1.450,7.630) {误解二：只要客观上不具备数额，就不定罪量刑};
\draw[line width=.4pt] (2.900,0.815) rectangle (7.300,6.330);
\node[align=left,text width=4.140cm,font={\linespread{1}\fontsize{10}{13.50}\selectfont}] at (5.100,3.572) {（1）只要客观上未达到“数额较大”，就不定罪。然而，盗窃罪的成立条件是，主观上有盗窃数额较大财物的故意，客观上有盗窃到数额较大财物的可能性、危险性。既遂条件是，盗窃到数额较大的财物};
\draw[line width=.4pt] (7.300,0.815) rectangle (16.000,6.330);
\node[align=left,text width=8.440cm,font={\linespread{1}\fontsize{10}{13.50}\selectfont}] at (11.650,3.572) {甲在商场盗窃到女士乙的提包，发现里面只有100元（不考虑提包本身的价值）。有意见认为，由于该财物未达到数额较大，因此甲不构成盗窃罪。然而，甲应成立盗窃罪，因为甲有盗窃到数额较大财物的可能性、危险性。甲构成犯罪未遂};
\draw[line width=.4pt] (2.900,6.330) rectangle (7.300,14.445);
\node[align=left,text width=4.140cm,font={\linespread{1}\fontsize{10}{13.50}\selectfont}] at (5.100,10.387) {（2）只要客观上未达到“数额（特别）巨大”，就不适用“数额（特别）巨大”的法定刑。然而，多数说认为，只要主观上有盗窃数额巨大财物的故意，客观上有盗窃到数额巨大财物的可能性，也即对该财物具有危险性，那么就应适用“数额巨大”的法定刑；若未盗窃到，再适用未遂规定};
\draw[line width=.4pt] (7.300,6.330) rectangle (16.000,14.445);
\node[align=left,text width=8.440cm,font={\linespread{1}\fontsize{10}{13.50}\selectfont}] at (11.650,10.387) {甲欲盗窃某珠宝店3万元的珠宝，手伸进玻璃柜（此时已经属于盗窃着手），左边是3万元的珠宝，右边是3千元的珠宝。甲认错对象，拿走3千元的珠宝。（1）多数说（司法解释）认为，甲应构成盗窃“数额巨大财物”的未遂，①同时，甲构成盗窃“数额较大财物”的既遂。甲的一个行为符合两个法定刑条件，想象竞合，择一重刑论处。理由是，盗窃“数额巨大财物”这种法定刑升格条件属于加重型构成要件，存在既遂与未遂问题。（2）少数说认为，由于甲未盗窃到数额巨大的财物，因此不能适用“数额巨大”档次的法定刑。理由是，“数额巨大”只有客观具备与否问题，没有既遂与未遂问题，属于单纯的量刑规则。对甲只能适用盗窃“数额较大财物”的既遂};
\end{tikzpicture}
